\documentclass{beamer}
\mode<presentation>
\usepackage{amsmath,amssymb,mathtools}
\usepackage{textcomp}
\usepackage{gensymb}
\usepackage{adjustbox}
\usepackage{subcaption}
\usepackage{enumitem}
\usepackage{multicol}
\usepackage{listings}
\usepackage{url}
\usepackage{graphicx} % <-- needed for images
\usepackage{xcolor}
\def\UrlBreaks{\do\/\do-}

\usetheme{Boadilla}
\usecolortheme{lily}
\setbeamertemplate{footline}{
	\leavevmode%
	\hbox{%
		\begin{beamercolorbox}[wd=\paperwidth,ht=2ex,dp=1ex,right]{author in head/foot}%
			\insertframenumber{} / \inserttotalframenumber\hspace*{2ex}
	\end{beamercolorbox}}%
	\vskip0pt%
}
\setbeamertemplate{navigation symbols}{}

\lstset{
	frame=single,
	breaklines=true,
	columns=fullflexible,
	basicstyle=\ttfamily\tiny    % tiny font so code fits
}

\numberwithin{equation}{section}

% ---- your macros ----
\newcommand{\conj}[1]{{#1}^c} % Custom command for conjugate
\providecommand{\nCr}[2]{\,^{#1}C_{#2}}
\providecommand{\nPr}[2]{\,^{#1}P_{#2}}
\providecommand{\mbf}{\mathbf}
\providecommand{\pr}[1]{\ensuremath{\Pr\left(#1\right)}}
\providecommand{\qfunc}[1]{\ensuremath{Q\left(#1\right)}}
\providecommand{\sbrak}[1]{\ensuremath{{}\left[#1\right]}}
\providecommand{\lsbrak}[1]{\ensuremath{{}\left[#1\right.}}
\providecommand{\rsbrak}[1]{\ensuremath{\left.#1\right]}}
\providecommand{\brak}[1]{\ensuremath{\left(#1\right)}}
\providecommand{\lbrak}[1]{\ensuremath{\left(#1\right.}}
\providecommand{\rbrak}[1]{\ensuremath{\left.#1\right)}}
\providecommand{\cbrak}[1]{\ensuremath{\left\{#1\right\}}}
\providecommand{\lcbrak}[1]{\ensuremath{\left\{#1\right.}}
\providecommand{\rcbrak}[1]{\ensuremath{\left.#1\right\}}}
\theoremstyle{remark}
\newtheorem{rem}{Remark}
\newcommand{\sgn}{\mathop{\mathrm{sgn}}}
\providecommand{\abs}[1]{\left\vert#1\right\vert}
\providecommand{\res}[1]{\Res\displaylimits_{#1}}
\providecommand{\norm}[1]{\lVert#1\rVert}
\providecommand{\mtx}[1]{\mathbf{#1}}
\providecommand{\mean}[1]{E\left[ #1 \right]}
\providecommand{\fourier}{\overset{\mathcal{F}}{ \rightleftharpoons}}
\providecommand{\system}{\overset{\mathcal{H}}{ \longleftrightarrow}}
\providecommand{\dec}[2]{\ensuremath{\overset{#1}{\underset{#2}{\gtrless}}}}
\newcommand{\myvec}[1]{\ensuremath{\begin{pmatrix}#1\end{pmatrix}}}
\newcommand{\mydet}[1]{\ensuremath{\begin{vmatrix}#1\end{vmatrix}}}

\newenvironment{amatrix}[1]{%
	\left[\begin{array}{@{}*{#1}{c}|c@{}}
	}{%
	\end{array}\right]
}

\newcommand{\myaugvec}[2]{\ensuremath{\begin{amatrix}{#1}#2\end{amatrix}}}
\let\vec\mathbf
% ---------------------

\title{12.367}
\author{EE25BTECH11042 - Nipun Dasari}

\begin{document}
	
	\begin{frame}
		\titlepage
	\end{frame}
	
	\begin{frame}{Question}
		The eigen-values of a real symmetric matrix are always
		\begin{enumerate}[label=\alph*)]
			\item positive
			\item imaginary
			\item real
			\item complex conjugate pairs
		\end{enumerate}
	\end{frame}
	
	\begin{frame}{Theoretical solution}
		Let $\vec{A}$ be a real symmetric matrix, meaning $\vec{A} = \vec{A}^\top$ and its entries are real. Let $\lambda$ be an eigenvalue and $\vec{v}$ its eigenvector. The eigenvalue equation is:
		\begin{align}
			\vec{A}\vec{v} = \lambda\vec{v}
		\end{align}
		Taking the complex conjugate of the equation:
		\begin{align}
			(\vec{A}\vec{v})^c = (\lambda\vec{v})^c \implies \vec{A}^c\vec{v}^c = \lambda^c\vec{v}^c
		\end{align}
		Since $\vec{A}$ is a real matrix, $\vec{A}^c = \vec{A}$. The equation becomes:
		\begin{align}
			\vec{A}\vec{v}^c = \lambda^c\vec{v}^c \label{eq:conj}
		\end{align}
	\end{frame}
	
	\begin{frame}{Theoretical solution}
		First, left-multiply the original eigenvalue equation by $(\vec{v}^c)^\top$:
		\begin{align}
			(\vec{v}^c)^\top\vec{A}\vec{v} = (\vec{v}^c)^\top(\lambda\vec{v}) = \lambda((\vec{v}^c)^\top\vec{v}) \label{eq:first}
		\end{align}
		Next, take the transpose of equation (\ref{eq:conj}):
		\begin{align}
			(\vec{A}\vec{v}^c)^\top = (\lambda^c\vec{v}^c)^\top \implies (\vec{v}^c)^\top\vec{A}^\top = \lambda^c(\vec{v}^c)^\top
		\end{align}
		Since $\vec{A}$ is symmetric, $\vec{A}^\top = \vec{A}$.
		\begin{align}
			(\vec{v}^c)^\top\vec{A} = \lambda^c(\vec{v}^c)^\top
		\end{align}
	\end{frame}
	
	\begin{frame}{Theoretical solution}
		Now, right-multiply this equation by $\vec{v}$:
		\begin{align}
			(\vec{v}^c)^\top\vec{A}\vec{v} = \lambda^c((\vec{v}^c)^\top\vec{v}) \label{eq:second}
		\end{align}
		By comparing equations (\ref{eq:first}) and (\ref{eq:second}), we see that:
		\begin{align}
			\lambda((\vec{v}^c)^\top\vec{v}) = \lambda^c((\vec{v}^c)^\top\vec{v}) \\
			(\lambda - \lambda^c)((\vec{v}^c)^\top\vec{v}) = 0
		\end{align}
	\end{frame}
	
	\begin{frame}{Theoretical solution}
		The term $(\vec{v}^c)^\top\vec{v}$ is the squared norm of $\vec{v}$:
		\begin{align}
			(\vec{v}^c)^\top\vec{v} = \sum_{i} v_i^c v_i = \sum_{i} |v_i|^2 = \norm{\vec{v}}^2
		\end{align}
		Since $\vec{v}$ is an eigenvector, it cannot be the zero vector, so its norm $\norm{\vec{v}}^2$ must be a positive real number.
		
		Given that $\norm{\vec{v}}^2 \neq 0$, the only way for the equation $(\lambda - \lambda^c)\norm{\vec{v}}^2 = 0$ to hold is if:
		\begin{align}
			\lambda - \lambda^c = 0 \implies \lambda = \lambda^c
		\end{align}
		A number is equal to its own conjugate only if its imaginary part is zero. Therefore, $\lambda$ must be real. The correct option is (c).
	\end{frame}
	
\end{document}